\documentclass[11pt]{article}

\usepackage[margin=1in]{geometry}
\usepackage{amsmath, amssymb}
\usepackage{tikz-cd}
\usepackage{lmodern}
\usepackage[T1]{fontenc}

\title{Tensor Algebra for the Encryption Algorithm:\\Extension to Image, Audio, Text, and Other Modalities}
\author{Alshicrypt Multimodal}
\date{\today}

\begin{document}
\maketitle

\section*{Objects}
\begin{itemize}
    \item $V$: tensor space for a modality
    \item $x \in V$: data tensor
    \item $T : V \to V$: tensor morphism / endomorphism
    \item $E : V \to V$: learned Encrypter
    \item $D : V \to V$: learned Decrypter
\end{itemize}

\section*{Tensor-Morphism Picture}
\[
\begin{tikzcd}[column sep=huge, row sep=large]
x \arrow[r, "T"] \arrow[d, "E"'] & T(x) \\
E(x) \arrow[ur, dashed, "\approx"'] &
\end{tikzcd}
\]

This expresses the intended learned relationship
\[
E(x) \approx T(x).
\]

\section*{Learned Reconstruction Picture}
\[
x \xrightarrow{\;E\;} E(x) \xrightarrow{\;D\;} D(E(x))
\]

with the intended reconstruction relation
\[
D(E(x)) \approx x.
\]

\section*{Interpretation}
The object $x$ is a full tensor, and the morphism $T$ acts on that tensor as a
whole. This is the tensor-algebra viewpoint: we treat the modality object
itself as a tensor in a tensor space $V$ and study endomorphisms
\[
T : V \to V.
\]

In the learned system, the Encrypter $E$ is trained to approximate the forward
tensor morphism, while the Decrypter $D$ is trained to map the transformed
tensor back toward the original.

\section*{Multimodal Examples}
\begin{itemize}
    \item For images, $x$ may be an RGBA tensor in $[0,1]^{H \times W \times 4}$.
    \item For audio, $x$ may be a waveform tensor or feature tensor in an
    analogous tensor space.
\end{itemize}

The current repository demonstrates the image case first, but the same
tensor-space framing is intended to extend to audio next and, more generally,
to text and other modalities that admit a tensor representation.

\end{document}
